\chapter{Dobór elementów obwodu mocy falownika}
\label{chap: Dobór elementów obwodu mocy falownika}

\section{Dobór tranzystorów MOSFET}
\label{sec: Dobór tranzystorów MOSFET}

\section{Dobór sterowników bramkowych}
\label{sec: Dobór sterowników bramkowych}

Jak wskazano w~podrozdziale~\ref{sec: Topologie jednofazowych falowników},
tranzystory górne mostka~H wymagają sterowników o~pływającym potencjale
odniesienia. Potencjał węzłów wyjściowych zmienia się przy każdym przełączeniu
między potencjałami biegunów N~i~P z~szybkością narastania rzędu kilku do
kilkudziesięciu woltów na nanosekundę. Aby załączyć tranzystor górny,
napięcie jego bramki musi być o~$U_\mathrm{GS(on)}$ wyższe od potencjału węzła
wyjściowego, czyli względem bieguna N wynosić nawet
$U_\mathrm{DC} + U_\mathrm{GS(on)}$. Wynikają z~tego trzy wymagania wobec
toru sterowania bramkowego:
\begin{itemize}
  \item \textbf{Zasilanie odniesione do źródła tranzystora.} Napięcie
    zasilające sterownik musi być podane względem źródła tranzystora górnego,
    a~nie względem masy układu~\cite[s.~29, 142]{mohan2012}. Stosuje się
    w~tym celu układ bootstrap albo izolowaną przetwornicę DC/DC
    z~transformatorem~\cite[s.~15]{instrukcja}, \cite[rozdz.~5, 7]{balogh2018}.
  \item \textbf{Przeniesienie sygnału sterującego.} Sygnał PWM z~mikrokontrolera,
    odniesiony do masy układu sterowania, trzeba przenieść do sterownika
    o~pływającym potencjale -- przez układ przesuwania poziomu
    (ang.~\textit{level shifter}) albo przez izolację galwaniczną (izolator
    cyfrowy, transoptor)~\cite[s.~29]{mohan2012}, \cite[rozdz.~5]{balogh2018}.
  \item \textbf{Odporność na szybkie zmiany potencjału.} Pojemności
    pasożytnicze między stroną pływającą a~masą układu sterowania (pojemność
    międzyuzwojeniowa transformatora przetwornicy, pojemność bariery
    izolatora) przewodzą przy każdym przełączeniu prąd
    $i = C\,\mathrm{d}u/\mathrm{d}t$. Przy pojemności \SI{10}{\pico\farad}
    i~szybkości narastania \SI{50}{\volt\per\nano\second} jest to impuls
    prądu o~wartości \SI{0,5}{\ampere}, który może zakłócić pracę sterownika
    lub mikrokontrolera. Elementy izolujące powinny więc mieć możliwie małą
    pojemność sprzęgającą.
\end{itemize}

Układ bootstrap jest prosty i~tani: kondensator zasilający sterownik górny
ładuje się przez diodę z~zasilania strony dolnej w~czasie, gdy przewodzi
tranzystor dolny tej samej gałęzi. Z~tego powodu nie pozwala on na długotrwałe
załączenie tranzystora górnego (współczynnik wypełnienia bliski 1), a~uzyskanie
ujemnego napięcia wyłączającego wymaga dodatkowych elementów~\cite[rozdz.~5]{balogh2018}.
Izolowane zasilanie każdego kanału jest droższe i~zajmuje więcej miejsca, ale
umożliwia pracę z~dowolnym współczynnikiem wypełnienia (również stałe
załączenie łącznika) oraz stosowanie ujemnego napięcia wyłączającego bramkę,
co zwiększa odporność na pasożytnicze załączenie tranzystora. Z~tych powodów
w~projektowanym stanowisku zastosowano izolowane zasilanie każdego z~czterech
kanałów sterowania bramkowego.


\section{Dobór kondensatorów magistrali DC}
\label{sec: Dobór kondensatorów magistrali DC}

\section{Analiza strat mocy}
\label{sec: Analiza strat mocy}

\section{Dobór radiatora}
\label{sec: Dobór radiatora}
